\documentclass[../../../include/open-logic-section]{subfiles}

\begin{document}

\olfileid{sth}{choice}{justifications}
\olsection{निवडीविषयीचे आंतरिक विचार}

मग व्यापक प्रश्न असा आहे की सुक्रमणाचे, निवडीचे किंवा सर्व संचांच्या आकारमानांच्या तुलनाक्षमतेचे---यांपैकी कोणतेही म्हणा---समर्थन करता येते का?

पुढे \emph{आंतरिक} समर्थनाचा एक प्रयत्न पाहू. \olref[z][story]{sec} मध्ये श्रेणीविषयी काही तत्त्वे मांडली होती. त्यांपैकी एक पुन्हा सांगण्याजोगे आहे:
\begin{enumerate}
	\item[] \stagesacc. कोणत्याही टप्पा $S$ साठी, आणि $S$ च्या \emph{आधी} तयार झालेल्या कोणत्याही संचांसाठी: नेमके ते संच घटक असलेला संच $S$ या टप्प्यावर तयार होतो. $S$ या टप्प्यावर दुसरे काहीही तयार होत नाही.
\end{enumerate}
खरे तर या तत्त्वाद्वारे किंवा यासारख्या तत्त्वाद्वारे निवडीच्या स्वयंसिद्धकाचे समर्थन करता येते, असे अनेक लेखकांनी सुचवले आहे. या पद्धतीचे थोडक्यात स्पष्टीकरण देऊ.

निवडीशी सममूल्य असलेले \emph{आणखी एक} तत्त्व देणाऱ्या साध्या निष्कर्षापासून सुरुवात करू:

\begin{thm}[$\ZF$ मध्ये]\ollabel{choiceset}
निवड पुढील तत्त्वाशी सममूल्य आहे. $A$ चे !!{element}s वियुक्त आणि अरिक्त असतील, तर प्रत्येक $x \in A$ साठी $C
\cap x$ एकघटकी असलेला $C$ अस्तित्वात आहे. (अशा $C$ ला $A$ साठीचा {निवडसंच} म्हणतो.)
\end{thm}

या निष्कर्षाचा पुरावा सरळ आहे आणि वाचकासाठी सराव म्हणून सोडला आहे.

\begin{prob}
\olref[sth][choice][justifications]{choiceset} सिद्ध करा. अडचण आल्यास \cite[pp.~242--3]{Potter2004} मध्ये पुरावा पाहता येईल.
\end{prob}

मुख्य मुद्दा असा की $A$ साठीच्या निवडसंचातून आधारकुटुंबाच्या युतीकरणाबाहेरील घटक वगळल्यावर तो $A$ साठीच्या निवडफलनाची प्रतिमा ठरतो. म्हणून निवडीचे समर्थन करण्यासाठी निवडसंचांच्या अस्तित्वाच्या भाषेतील तिच्या सममूल्य मांडणीचे समर्थन करण्याचा प्रयत्न करता येईल. आता नेमके तेच करू.

$A$ चे !!{element}s वियुक्त आणि अरिक्त असू द्या. \stageshier{} नुसार (पाहा \olref[z][story]{sec}) $A$ कोणत्या तरी टप्पा $S$ वर तयार होतो. $\bigcup A$ चे सर्व !!{element}s $S$ टप्प्याच्या आधी उपलब्ध असतात. आता \stagesacc{} नुसार, $S$ च्या आधी तयार झालेल्या \emph{कोणत्याही} संचांचे नेमके तेच घटक असलेला संच तयार होतो. दुसऱ्या शब्दांत: आधी उपलब्ध झालेल्या संचांचे प्रत्येक \emph{शक्य} संकलन $S$ वर अस्तित्वात असते. पण $A$ साठीच्या निवडसंचात एकत्र करता येतील अशा वस्तू निवडणे नक्कीच \emph{शक्य} आहे, असा या समर्थनाचा दावा आहे; ते $\bigcup A$ च्या अगदी विशिष्ट उपसंचातच येतील. म्हणून आवश्यक असा एखादा निवडसंच अस्तित्वात आहे, असे या युक्तिवादातून म्हटले जाते.

आंतरिक आधारावर निवडीचे समर्थन करण्याचा हा \emph{अतिशय} थोडक्यात मांडलेला प्रयत्न आहे. ही कल्पना पुढे नेण्यासाठी पॉटर यांनी (\citeyear[\S14.8]{Potter2004}) केलेली सुबक मांडणी वाचावी.

\end{document}
